% Transcription — fonds Alexandre Grothendieck, shelfmark 134-2, batch 1 (pages 1-20).
% Established from the facsimile archives/batches/134-2/batch-01.pdf, per the
% skill transcribe-grothendieck and the special rules for folder 134-2 (#44).
%
% Pass: Opus 5.5 (claude-opus-5-5), 2026-10-03 — first pass, unchecked against the pages by a human.
%
% What the batch holds. Folder 134-2 opens Pursuing Stacks.
%   Page 1 — a cover in his pencil, « The Modelizing Story »: no \page{}, his
%     words become the section heading.
%   Pages 2-12 — the provisional plan and tables of contents, typed and heavily
%     revised by hand (pages 8 and 10 wholly manuscript): p. 2 the plan in
%     parts I-X with a dependency sketch; p. 3 his p. « i », first table of
%     chapter I, the Appendix (App. 1-18) of the letters to Breen, and the
%     start of II; pp. 4-7 the later table (« PURSUING STACKS / First episode:
%     THE MODELIZING STORY ») §§ 1-66; pp. 8-12 §§ 67-140 (parts IV, VI, VII).
%     All transcribed: they are his tables, not the text of the sections.
%   Page 13 — manuscript dedication (to his former students, struck; to his
%     brother Dodek Shapiro). Transcribed.
%   Pages 14-20 — the letter to Daniel Quillen, « Les Aumettes 19.2.1983 »,
%     his pp. 1-6, Pursuing Stacks §§ 1-4 (p. 20 opens « 20.2. » and § 4). Not
%     transcribed (printed by Maltsiniotis, SMF Documents mathématiques 20,
%     2022); each page carries \page{} and one note. The letter runs on past
%     the end of the batch.
%
% Letters in this batch:
%   pp. 14-20 — to Daniel Quillen, « Les Aumettes 19.2.1983 » (cont. 20.2),
%     English, his own: skipped as Pursuing Stacks text (§§ 1-4), per #44.
%   No letter to Breen and no third-party letter in pages 1-20; the tables on
%     pp. 3-4 only announce the Appendix « Three letters to Larry Breen ».
%
% Unclear whether Pursuing Stacks text: p. 13 (the dedication — not a numbered
% section, so transcribed, but the published edition may print it); p. 14,
% whose pasted head slip « Chapter I : TAKE OFF / A letter to …. / (followed
% by an Appendix: Three letters to Larry Breen) » sits above § 1 — treated with
% the letter and not transcribed. The tables (pp. 2-12) are his working tables
% and differ from any printed one; transcribed.
%
% The dating is the inventory's own for the shelfmark (catalogue, group 134-1-8).
\documentclass[11pt,a4paper]{article}
\input{../preamble/grothendieck}

\folder{134-2}
\batch{1}
\pages{1}{20}
\foldertitle{[Chapitre I : Take off (pages 1 à 65) et table des matières provisoire] : tapuscrits annotés (19/02-22/02/1983), note manuscrite (s.d.), copies de lettre (1975, s.d.).}
\dating{1975-[1983]}
\watermark{Édition de démonstration}

\begin{document}

\section{The Modelizing Story}
\note{titre écrit au crayon, de sa main, sur le feuillet de couverture (p. 1), qui ne porte rien d'autre.}

\page{2}
\begin{quote}
THE MODELIZING STORY

(Histoire de Modèles)
\end{quote}

\begin{itemize}
  \item[I] Take-off
  \item[II] Test categories \struck{and test functors}
  \item[III] Homotopy structures \struck{and contractibility structures}
  \item[IV] Asphericity structures \struck{and canonical modelizers}
  \item[V] \struck{Homology and cohomology (abelianization of homotopy types)} \add{Abelianization}
  \item[VI] \add{Schematization}
  \item[VIII] Homotopy properties of (Cat) and $A^{\wedge}$ (closed model structures)
  \item[IX] Derivators
  \item[X] Back to topoi
\end{itemize}
\note{table des parties, dactylographiée et remaniée à la main. Les numéros sont corrigés à la main : « VII » devient « VIII » (Homotopy properties…), « VIII » devient « IX » (Derivators), « VIII » biffé puis « X » (Back to topoi) ; sous « V » un « VI » biffé. Des traits de crayon renvoient « Test categories », « Homotopy structures » et « Asphericity structures » vers une même accolade marquée « I » en marge droite ; au crayon, à peine lisible, « \uncertain{homotopy types} » repris deux fois dans les cases ajoutées autour de V et VI.}
\marginal{\add{VII} Abelianization (II)}
\note{écrit verticalement dans la marge gauche, en face des parties V à VIII : une partie VII « Abelianization (II) » s'ajoute ainsi au plan.}

\note{sous la liste, au crayon puis à l'encre, un schéma des dépendances entre parties :}

\begin{tikzcd}[column sep=small, row sep=small]
\mathrm{I} & \mathrm{II} \arrow[r, Rightarrow] & \mathrm{III} \arrow[r, Rightarrow] & \mathrm{IV} \arrow[dd, Rightarrow] & \mathrm{V} \arrow[d, Rightarrow] \arrow[dr] & \\
 & & & & \mathrm{VII} & \mathrm{VI} \\
 & & & \mathrm{VIII} \arrow[d, Rightarrow] & & \\
 & & & \mathrm{IX} & &
\end{tikzcd}

\note{deux flèches biffées : l'une de V vers un « VI » biffé placé à droite de V sur la première ligne, l'autre de V vers le bas à gauche, en direction de VIII.}

\page{3}
\note{p. i de l'auteur.}
\begin{quote}
TABLE OF CONTENTS of The Modilizing Story (1).
\end{quote}

\marginal{NB Les chiffres soulignés 5, 9, 13, App.\,3, App.\,6, App.\,11, App.\,12 etc sont à imprimer en caractères gras}

I \emph{Take-off} ( A letter to … )
\begin{itemize}
  \item[1.] The importance of innocence.
  \item[2.] A short look into purgatory.
  \item[3.] ``Fundamental $\infty$-groupoids'' as objects of a ``model category'' ?
  \item[4.] A bit of ordering in the mess of ``higher order structures''.
  \item[5.] Jumping over the abyss !
  \item[6.] The topological model; hemispheres building up the (tentative) ``universal $\infty$-(co)groupoid''.
  \item[7.] Gluing hemispheres: the ``standard'' amalgamations.
  \item[8.] Description of the universal \emph{primitive} structure.
  \item[9.] The main inductive step: just add coherence arrows ! And the abridged story of an (unescapable and irrelevant) ambiguity …
  \item[10.] Cutting down redundancies -- or: ``l'embarras du choix''.
  \item[11.] Back to the topological model. The canonical functor from spaces to ``$\infty$-groupoids''.
  \item[12.] About replacing spaces by objects of a ``model category''.
  \item[13.] An urgent reflection on proper names: ``stacks'' and ``coherators''.
\end{itemize}
\marginal{espacer deux fois}
\note{consignes de mise en page pour la frappe : « espacer deux fois » avant l'Appendix, et « espacer » avant App.\,6 et avant App.\,9.}

\emph{Appendix} : Three letters to Larry Breen
\begin{itemize}
  \item[(App.1)] Exemples de 2-catégories de Picard.
  \item[(App.2)] \struck{Les n-groupoi} Théorème de Lefschetz (faible) en termes de \add{(}Champs.\add{)}
  \item[(App.3)] Les n-groupoides comme types d'homotopie tronqués.
  \item[(App.4)] Relativisation sur un topos.
  \item[(App.5)] Ingrédients principaux vers une ``Algèbre Topologique''.
  \item[(App.6)] Champs essentiellement localement constants et types d'homotopie \add{(}pie\ill{}
  \note{la fin de la ligne déborde le bord de la feuille numérisée.}
  \item[(App.7)] Intégration de champs et cohomologie.
  \item[(App.8)] Les trois approches vers la cohomologie d'un topos.
  \item[(App.9)] \struck{Refle} Question about a non-abelian Dold-Kan theorem.
  \item[(App.10)] The ``six operations'' and homology.
  \item[(App.11)] Complex of generalized jacobians.
  \item[(App.12)] Global ``geometric'' class field theory as a cohomological duality formula. Serre duality and the ``Lang trick''.
  \item[(App.13)] Case of local ``geometric'' class field theory.
  \item[(App.14)] Comments on fppf cohomology \struck{and} \add{versus crystallin cohomology.}
  \item[(App.15)] The homotopical trinity: \add{(}``spaces'', ``$\infty$-groupoids'', \add{(}essentially locally constant) ``coefficient categories''…
  \item[(App.16)] The ``six operations'' for P-constructible sheaves (for a given equisingular stratification P).
  \item[(App.17)] Relation with ``indexed homotopy types'' and with ``dévissage'' of stratified spaces. The ``fine'' homotopy type of a (tame) space or of a scheme.
  \item[(App.18)] How far do ``essentially constant'' \emph{abelian} coefficients determine a homotopy type ?
\end{itemize}

II \emph{Test Categories}
\begin{itemize}
  \item[14.] The unnoticed failure (of the foundations of topology).
  \item[15.] Overall review on standard descriptions of homotopy types.
  \item[16.] Stacks over topoi as the unifying concept for homotopical and cohomological algebra.
  \item[16'.] Categories as models for homotopy types. First glimpse upon an ``impressive bunch'' (of modelizers).
  \item[17.] The Artin-Mazur cohomological criterion for weak equivalence.
  \item[18.] Corrections and comments to letter. Benabou's lonely approach.
\end{itemize}
\note{« 14 » est surchargé à la main d'un chiffre peu lisible, et « 16' » porte un accent ajouté à la main : la numérotation des §§ 14--18 est en cours de remaniement (cf. p. 5).}

\page{4}
\note{deuxième état de la table, plus tardif que celui de la p. 3, avec titre général.}
\begin{quote}
\add{PURSU}ING STACKS

( A la \add{P}oursuite des Champs )

First episode : THE MODELIZING STORY

(Histoire de Modèles )
\end{quote}
\note{le début du titre est récrit à la main en surcharge des premières lettres tapées, illisibles dessous ; « (Histoire de Modèles) » est entouré d'un trait de crayon.}

\emph{Contents}

I \emph{The Take-off} (a letter to Daniel Quillen)
\begin{itemize}
  \item[1.] The importance of innocence.
  \item[2.] A \struck{quick} \add{short} look into purgatory
  \item[3.] ``Fundamental $\infty$-groupoids'' as objects of a ``model category'' ?
  \item[4.] A bit of ordering in the mess of ``higher order structures''.
  \item[(5.)] Jumping over the abyss !
  \item[6.] The topological model: hemispheres building up the (tentative) ``universal $\infty$-(co)groupoid''.
  \item[7.] Gluing hemispheres: the ``standard'' amalgamations.
  \item[8.] Description of the universal \emph{primitive} structure.
  \item[(9.)] The main inductive step: just add coherence arrows ! \struck{The story of the irrelevant} \add{And the} abridged story \struck{ambiguity} of an (unescapable and irrelevant) ambiguity …
  \item[10.] Cutting down redundancies -- or: ``l'embarras du choix''.
  \item[11.] Returning to the topological model\struck{: getting} \add{.} The \struck{geometric realization for ($\infty$-groupoids)} canonical functor from spaces to ``$\infty$-groupoids''.
  \item[12.] About replacing spaces by objects of a ``model category''.
  \item[(13.)] \struck{Reflexion} An urgent reflection on proper names: ``Stacks'' and ``coherators''.
\end{itemize}
\note{les biffures de cette page sont en partie faites à la machine (x frappés) et les corrections retapées au-dessus de la ligne ; les numéros 5, 9 et 13 sont cerclés à l'encre. Le mot biffé au § 13 est peu lisible (\uncertain{Reflexion}).}

\struck{\ill{}} II \add{\emph{Test categories and test functors.}}
\begin{itemize}
  \item[14.] \struck{Homotopical algebra versus ``topology''.} The \struck{essential inadequacy} unnoticed \struck{breakdown} \add{failure} (of the present foundations of topology for expressing topological intuition).
  \item[15.] Overall review on \add{(}standard\add{)} descriptions of homotopy types. \struck{by models.}
  \item[16.] \struck{The main motivation: stacks over topoi.} Stacks \add{over topoi,} as the unifying concept for homotopical and cohomological algebra.
  \item[16 bis.] Categories as models for homotopy types. First glimpse upon an ``impressive bunch'' \add{(of modelizers).}
\end{itemize}
\note{à la fin du § 14, « essential inadequacy » (\uncertain{essential}) est biffé à la machine, « unnoticed breakdown » tapé au-dessus, puis « breakdown » biffé à la main et remplacé par « failure ». Au § 16 bis, le « 5 » est surchargé et « bis » ajouté à la main.}

\page{5}
\begin{itemize}
  \item[17.] The Artin-Mazur cohomological criterion for weak equivalence.
  \item[18.] Corrections and comments to \add{l}etter. Benabou's lonely approach.
  \item[19.] \struck{Starting} \add{Beginning of} a provisional itinerary \add{(}through stacks\add{)}.
  \item[20.] Are model categories sites ?
  \item[21.] Further glimpse upon the ``bunch'' of possible model categories\add{;} \struck{and} \add{and upon} relation between n-complex\add{es} and n-stacks.
  \item[22.] Ordered sets as models for homotopy types.
  \item[(23.)] Getting a basic functor $M\to(\mathrm{Hot})$ \struck{\ill{}} \add{from a site} structure $M$ (\struck{the} faltering beginning of a systematic reflection).
  \item[24.] A bunch of topologies on (Cat).
  \item[25.] A tentative equivalence relation for topologies.
  \item[26.] The dawn of \struck{\ill{}} \add{test categories} and test functors …
  \item[27.] Digression on ``geometric realization'' functors.
  \item[(28.)] The ``inspiring assumption''. Modelizers.
  \item[29.] The basic modelizer (Cat). Provisional definition of test categories and elementary modelizers.
  \item[30.] Starting \struck{with} the ``asphericity game''.
  \item[(31.)] The end of the thin air conjecturing: \struck{characterization of} \add{a criterion for} test categories.
  \item[32.] Provisional program of work.
  \item[(33.)] Necessity of conditions T\,1 to T\,3\add{.} \struck{And translation in terms of elementary modelizers.}
  \item[34.] Examples of test categories.
  \item[35.] The notion of a modelizing topos. Need \struck{of} \add{for} revising the \add{Č}ech-Verdier-Artin-Mazur construction.
  \item[36.] Characterization of \add{(a particular type of)} test functors with values in (Cat).
  \item[(37.)] The ``asphericity story'' told anew -- the ``key result'' \add{(on test functors)}. \add{[Homotopy intervals]}
  \item[38.] Asphericity story retold (contd): generalized nerve functors.
  \item[39.] Returning upon terminology: strict test categories, and strict modelizers.
  \item[40.] Digression on cartesian products of weak equivalences in (Cat)\add{;} weak equivalences \emph{relative} to a given base object.
  \item[41.] Role of the ``inspiring assumption'', and of saturation conditions on ``weak equivalences''.
\end{itemize}
\note{numéros 23, 28, 31, 33 et 37 cerclés ; une flèche au crayon relie 31 à 33. Des traits de crayon en marge gauche groupent les §§ en blocs (après 19, 23, 27, 29, 31, 34, 36, 37, 41). « [Homotopy intervals] » est au crayon, en marge droite du § 37, lecture \uncertain{Homotopy intervals}. Au § 40, le « 4 » en tête de la seconde ligne est un chiffre biffé.}

\page{6}
\begin{itemize}
  \item[42.] Terminology\struck{x} revised (model preserving functors). Submodelizers of the basic modelizer (Cat).
  \item[43.] The category $\Delta^{f}$ of simplices without degeneracies as a weak test category -- or ``face complexes'' as models for homotopy types.
  \item[44.] Overall review of the basic notions:
  a) weak test categories ;
  b) test categories and local test categories ;
  c) strict test categories ;
  d) weak test functors and test functors \add{(}with values in (Cat)).
\end{itemize}

III \struck{Grinding through abstract homotopy notions,} \add{(Homotopy structures,} \struck{my way towards} \add{contractibility structures} \struck{canonical modelizers}\add{)}
\note{intitulé de partie récrit à la main sur une ligne dactylographiée biffée ; les mots « contractibility » et « structures » sont écrits au-dessus de la ligne.}

\begin{itemize}
  \item[45.] It's burning again ! Review of some ``recurring striking features'' of modelizers\add{,} and \add{of} standard \struck{\ill{}} \add{modelizing} functors.
  \item[(46.)] Test functors with values in any modelizer: an observation, and an \add{inspiring} ``silly question''.
  \item[47.] An approach for handling (Cat)-valued test functors\add{,} \add{(and promise of a ``key result'' revised.)} The significance of contractibility.
  \item[48.] A journey through \add{abstract} homotopy notions\note{un trait de plume traverse « homo- » ; « abstract » est inséré au-dessus.}\add{,} in terms of a set W of ``weak equivalences''.
  \item[49.] Contractible objects. Multiplicative intervals.
  \item[(50.)]\note{numéro « 51 » corrigé à la main en « 50 », cerclé.} The four basic ``pure'' homotopy notions\struck{ (\uncertain{foresight of idyllic})} and their interplay: \struck{a sug\ill{}} \struck{``canonical modelizers''} \struck{(with variations …)}.
  \note{intitulé surchargé à la machine : les passages biffés le sont par des x dactylographiés ou d'un trait à la main.}

  A) Homotopy relation between maps.

  B) Homotopisms, and homotopism structures.

  C) Homotopy interval structures.

  D) Contractibility structures.

  E) Generating sets of homotopy intervals. Two standard ways of getting multiplicative intervals. Contractibility of $\underline{\mathrm{Hom}}(X,Y)$'s.

  F) The canonical homotopy structure: preliminaries on $\pi_{0}$.
  \item[(51.)]\note{numéro « 52 » corrigé à la main en « 51 », cerclé.} Reflection on some main impressions. The foresight of an ``idyllic picture'' (of would-be ``canonical modelizers'').
  \note{une flèche à la main relie le § 51 au § 50 et en inverse l'ordre.}
  \item[52.] Inaccuracies rectified.
  \item[53.] Compatibility of a functor $u\colon M\to N$ with a homotopy structure on $M$.
  \item[54.] Compatibility of a homotopy structure $h$ with a set $W$ of ``weak equivalences''. The homotopy structure $h_{W}$.
  \item[55.] Maps between homotopy structures.
\end{itemize}

\page{7}
\begin{itemize}
  \item[56.] Another glimpse upon canonical modelizers. Provisional working plan -- and recollection of some questions.
  \item[57.] Relation of homotopy structures to 0-connectedness and $\pi_{0}$. The canonical homotopy structure $h_{M}$ of a category $M$.
  \item[58.] Case of totally 0-connected category $M$. The category $\overline{(\mathrm{Cat})}$ of (small) categories and homotopy classes of functors.
  \item[59.] Case of the ``next best'' modelizer (Spaces) -- and need of introducing the $\pi_{0}$-functor as an extra structure on a would-be modelizer $M$.
  \item[60.] Case of \add{a} strictly totally aspheric topos. A timid start on axiomatizing the set $\underline{W}$ of weak equivalences in (Cat).
  \item[61.] Remembering about the \add{(promised)} ``key result'' at last!
  \item[(62.)] An \add{(embarassing)} case of hasty over-axiomatization. The unexpected richess …
  \item[63.] Review of terminology (provisional).
  \item[64.] Review of properties of the ``basic localizer'' $\underline{W}_{(\mathrm{Cat})}$.
  \item[65.] Still another review of the test notions \add{(relative to given basic localizer)}:
  A) Total \add{$\underline{W}$-}asphericity.
  B) Weak $\underline{W}$-test categories.
  C) $\underline{W}$-test categories.
  D) Strict $\underline{W}$-test categories.
  E) Weak $\underline{W}$-test functors, and $\underline{W}$-test functors. The ``key result'' in the long last ! (first version.) \add{cf § 79 for ``final shape''…}
  F) $\underline{W}$-test functors $A\to(\mathrm{Cat})$ of strict $\underline{W}$-test categories.
  \item[66.] Revising (and fixing ?) terminology again.
\end{itemize}
\marginal{p. 168 --}
\note{en marge gauche du § 65, à l'encre, « p. 168 -- » suivi d'un mot ou d'un chiffre peu lisible (\uncertain{199}) : renvoi à la pagination du tapuscrit.}

\struck{IV ASPHERICITY STRUCTURES AND CANONICAL MODELIZERS.}
\begin{itemize}
  \item[\struck{67.}] \struck{Setting out for the asphericity game again: variance of the category $(\mathrm{Hot}_{A})$, for arbitrary small category $A$ and aspheric functors.}
  \item[\struck{68.}] \struck{Digression on a ``New Continent''.}
  \item[\struck{69.}] \struck{Digression on six weeks' scratchwork (1): derivators, and integration of homotopy types.}
  \item[\struck{70.}] \struck{Digression on scratchwork (2): cohomological properties of maps in (Cat), and} \add{\struck{in}} \struck{\ill{} $A^{\wedge}$.} \struck{Kan fibrations \ill{} a duality} \struck{Foreboding of} \struck{\ill{}} \struck{Does any topos admit a ``dual'' topos ?} \struck{Kan fibrations rehabilitated !}
  \item[\struck{71.}] \struck{Working program, and rambling questions (notably on group objects as models, and on the Dold-Puppe theorem…).}
\end{itemize}
\note{tout le bas de la page, de l'intitulé IV au § 71, est barré de grandes croix : la suite de la table est reprise plus loin. Les §§ 68 à 71 sont écrits à la main ; au § 70 plusieurs reprises biffées se chevauchent, et leur ordre de lecture est incertain.}

\page{8}
\note{page manuscrite : la table reprend au chapitre IV, au propre.}

IV \struck{\ill{}} \emph{Asphericity structures} \struck{and canonical modelizers}.
\begin{itemize}
  \item[67.] Back to the asphericity game: the categories $\mathrm{Hot}_{A}$.
  \item[68.] Digression on a ``new continent''.
  \item[69.] Digression on six weeks' scratchwork: \emph{derivators}, and integration of homotopy types.
  \item[70.] Digression on scratchwork (2): cohomological properties of maps in (Cat) and in $A^{\wedge}$. Does any topos admit a ``dual'' topos? Kan fibrations rehabilitated.
  \item[71.] Working program and rambling questions (group objects as models, Dold-Puppe theorem…).
  \item[(72.)] Back to asphericity: criteria for a map in (Cat).
  \item[73.] Asphericity criteria (contd).
  \item[74.] Application to products of test categories.
  \item[(75.)] Asphericity structures: a bunch of useful tautologies.
  \item[76.] Examples. Totally aspheric asphericity structures.
  \item[(77.)] The canonical functor $\mathrm{Hot}_{M}\to(\mathrm{Hot})$.
  \item[78.] Test functors and modelizing asphericity structures: the outcome (at last!) of an early ``observation''.
  \item[(79.)] Asphericity structure generated by a contractibility structure: the final shape of the ``awkward main result'' on test functors.
  \item[80.] Reminders and questions around canonical modelizers.
\end{itemize}
\note{numéros 72, 75, 77 et 79 cerclés au crayon ; traits de crayon en marge gauche après 67, 68, 70, 71, 72, 74, 77 et 78.}

\page{9}
\begin{itemize}
  \item[(81.)] Contractibility as the common expression of homotopy, asphericity and 0-connectedness notions. (An overall review of the notions met with so far.)
  \item[82.] Proof of injectivity of $\alpha\colon \mathrm{Contr}(M)\hookrightarrow \underline{W}\text{-}\mathrm{Asph}(M)$. Application to $\underline{\mathrm{Hom}}$ objects and to products of aspheric functors $A\to M$.
  \item[83.] Tautologies on $\mathrm{Im}\,\alpha$, and related questions.
  \item[84.] A silly (provisional) answer to the ``silly question'' \add{(cf section 90)} -- and the new perplexity $f_{!}(M_{\mathrm{as}})\subset M'_{\mathrm{as}}$ ?
  \item[85.] Digression on left exactness properties of $f_{!}$ functors, application to the inclusion $i\colon \Delta\hookrightarrow(\mathrm{Cat})$.
  \item[86.] Bimorphisms of contractibility structures as the (final ?) answer. Does the notion of a map of asphericity structures exist ?
  \item[87.] Comments on Thomason's paper on closed model structure of (Cat).
  \item[88.] Review of pending questions and topics (questions 1) to 5), including characterizing canonical modelizers).
  \item[89.] Digression (cont\struck{\ill{}}\add{d}) on left exactness properties of $f_{!}$ functors.
  \item[90.] Review of questions (contd): 6) Existence of test functors and related questions. Digression on strictly generating subcategories.
  \item[91.] Review of questions (contd.): 7) Homotopy types of finite type, 8) test categories with boundary operations, 9) miscellaneous.
  \item[92.] Short range working program, and an afterthought on abelianization of homotopy types: a handful of questions around the Whitehead and \add{Kan-}Dold-Puppe theorems. \add{(The unsuspected, awkward start of a systematic reflection on \struck{\ill{}} …)}
  \item[93.] The afterthought continued: abelianizators, and ``standard'' abelianizators for categories with boundary operators.
  \item[94.] Afterthought (contd): retrospective on the ``De Rham complex with divided powers'' and on some wishful thinking about linearization of homotopy types and \add{(arbitrary)} ground-ring extension in homotopy types.
  \item[95.] Contractors.
  \item[96.] ``Vertical'' and ``horizontal'' topoi … (afterthoughts on terminology).
  \item[97.] ``Projective'' topoi. Morphisms and bimorphisms of contractors.
  \item[98.] Sketch of proof of $\Delta^{f}$ being a weak test category -- and perplexities about it's being aspheric !
  \item[99.] Afterthoughts on abelianization IV: Integrators.
\end{itemize}
\note{les §§ 87 à 91 sont encadrés au crayon, et la limite supérieure du cadre tracée par un trait ondulé entre 86 et 87. Dans la marge gauche du cadre, à l'encre, en oblique et biffé : « V Homology and cohomology / abelianization / homotopy t\supplied{ypes} » ; un « V » noirci par deux taches d'encre. Les §§ 95 à 98 sont réunis par un crochet. Au § 90, le « 0 » surcharge un autre chiffre.}

\page{10}
\note{page manuscrite, suite de la table.}
\begin{itemize}
  \item[100.] Homology versus cohomology
  \item[101.] Abelian integration $\mathrm{L}f_{!}^{\mathrm{ab}}$ (versus abelian cointegration $\mathrm{R}f_{*}$)
  \item[(102.)] Integrators are abelianizators.
  \item[103.] Integrators versus cointegrators
  \item[104.] Overall review on abelianization (1): Case of \add{(pseudo-topoi.)}
  \item[105.] Review (2): duality equivalences for ``algebraic'' topoi\add{.} \struck{and abelian (topoi.}
  \item[106.] Review (3): A formulaire for the basic integration and cointegration operations $*$ and $\mathrm{Hom}$.
  \item[107.] Review (4): Case of general ground ring $k$.
  \item[108.] Review (5): Homology and cohomology (absolute case)
  \item[(109.)] Review (6): A further step in linearization: coalgebra structures $P\to P\otimes_{k}P$ in (Cat).
\end{itemize}
\note{au § 104, « pseudo-topoi » est écrit au-dessus de la ligne du § 105 et rattaché par un trait à « Case of » ; la place de l'insertion est \uncertain{celle-ci}.}

\marginal{VI Schematization.}
\begin{itemize}
  \item[(110.)] More wishful thinking on ``schematization'' of homotopy types.
  \item[(111.)] Complexes of ``unipotent bundles'' as models, and ``schematic'' linearization.
  \item[112.] Postnikov dévissage and Kan condition for schematic complexes.
  \item[113.] ``Soft'' versus ``hard'' Postnikov dévissage; $\pi_{1}$ as a group scheme.
  \item[114.] Outline of a program.
  \item[115.] $L(X)$ as the pro-quasicoherent substitute for $\underline{O}_{k}^{(X)}$.
  \item[116.] The need and the drive.
  \item[117.] ``Schematic'' versus ``formal'' homology and cohomology invariants.
  \item[(118.)] The homotopy groups $\pi_{i}$ as derived functors of the ``Lie functor''. Lack of satisfactory models for $S^{2}$ and $S^{3}$.
\end{itemize}
… (à suivre)
\note{un trait sépare les §§ 109 et 110 ; l'intitulé « VI Schematization » est en marge gauche, en oblique, face aux §§ 110--111.}

\page{11}
\begin{itemize}
  \item[(119.)] Breakdown of an idyllic picture -- and a tentative next best ``binomial'' version of the ``comparison theorem'' for schematic versus discrete linearization.
  \item[120.] Digression on the Lazare ``analyzers'' for ``binomial'' commutative algebra, $\lambda$-commutative algebra etc.
  \item[121.] The basic pair of adjoint functors $\mathrm{Hotab}_{0} \underset{\widetilde{K}}{\overset{\widetilde{LH}_{\bullet}}{\leftrightarrows}} \mathrm{Hot}_{0}$.
  \item[122.] Rambling reflections on $\widetilde{LH}_{\bullet}$, Postnikov invariants, $S(H,n)$'s -- and on the non-existence of a ``total homotopy''-object $L\pi_{\bullet}(X)$ \struck{for the discrete set-up.} \add{for ordinary homotopy types.}
  \item[123.] The hypothetical complexes ${}^{n}\Pi_{\bullet}=L\pi_{\bullet}(S^{n})$, and comments on homotopy groups of spheres.
  \item[(124.)] Outline of a program (second version): an autodual formulaire for the basic ``four functors'' $\widetilde{LH}_{\bullet}$, $L\pi_{\bullet}$, $\widetilde{K}$, $\widetilde{S}$.
  \item[125.] Digression on Baues' cofibration and fibration categories, and on weak equivalences alone as the basic data for a model category.
  \item[126.] Dissymetry of $\widetilde{LH}_{\bullet}$ versus $L\pi_{\bullet}$ (no filtration dual to Postnikov dévissage).
  \item[127.] Schematic homotopy types and Illusie's derived category for a topos.
  \item[128.] Looking for the right notion of ``bundles''; V-bundles versus W-bundles.
  \item[129.] Quasi-coherent homological quasi-isomorphisms, versus weak equivalences.
  \item[130.] A case for non-connected bundles.
  \item[(131.)] Tentative description of the spherical functor $\widetilde{S}$, and ``infinitesimal'' extension of the basic notion of ``bundles''.
  \item[(132.)] A crazy, tentative, wrong-quadrant (bi)complex for the homotopy groups of a sphere.
  \item[133.] Birth of Suleyman.
\end{itemize}
\note{numéros 119, 124, 131 et 132 cerclés ; un trait sépare les §§ 132 et 133. Les fins de ligne des §§ 127--129 (« topos. », « bundles. », « valences. ») sont retapées au-dessus de la ligne. Dans le tapuscrit, le point de $\widetilde{LH}.$, $L\pi.$ est le point bas d'indice gradué.}

\page{12}
VII \emph{Linearization of homotopy types} \struck{(II)} (2)
\begin{itemize}
  \item[133.] Birth of Suleyman.
  \item[134.] Back to linearization in the modelizer (Cat): another handful of questions around Kan-Dold-Puppe.
  \item[135.] Proof of ``integrators are abelianizators'' (block against homology receding ?).
  \item[136.] Preliminary perplexities about a full-fledged ``six operations duality formalism'' within (Cat).
  \item[137.] Looking for the relevant notions of properness and smoothness for maps in (Cat). Case of ordered sets as a paradigm for cohomology theory of \struck{\ill{}} canonically stratified spaces.
  \item[138.] Niceties and oddities: $\mathrm{R}f_{!}$ commutes to colocalization, not localization.
  \item[139.] Retrospective on ponderings on abelianization, and on coalgebra structures in (Cat).
  \item[140.] The meal and the guest.
\end{itemize}
… \add{(to be completed)}

\page{13}
\note{feuillet de dédicace, manuscrit. La première dédicace est barrée de traits obliques au crayon.}
\begin{quote}
\struck{À ceux qui furent mes élèves -- à qui j'ai donné du meilleur de moi-même, et aussi du pire …}

\struck{et} \struck{À} \add{À} mon frère, Dodek Shapiro, que j'aurais voulu rencontrer un jour, alors que les seuls signes tangibles \add{(que j'aie)} de sa vie sont quelques lettres et photos d'enfant, anciennes de près d'un demi-siècle, \struck{\ill{}} datées d'Ulianovsk en Sibérie, \struck{\ill{}} adressées à \struck{notre} \add{son} \struck{\uncertain{maman}} père, Sascha \struck{\ill{}} (qui est aussi le mien) …
\end{quote}
\note{« et » est biffé et le « À » repris en majuscule ; le mot biffé sous « Sibérie » est lu \uncertain{maman}, puis noirci.}

\page{14}
\note{Pursuing Stacks, § 1 (« Les Aumettes 19.2.1983 », lettre à Daniel Quillen, p. 1 de l'auteur, avec le bandeau collé « Chapter I : TAKE OFF ») ; imprimé dans Maltsiniotis (éd.), À la poursuite des champs, vol. I, SMF 2022. Non transcrit.}

\page{15}
\note{Pursuing Stacks, § 2 (lettre à Daniel Quillen, suite, p. 2 de l'auteur) ; imprimé dans Maltsiniotis (éd.), À la poursuite des champs, vol. I, SMF 2022. Non transcrit.}

\page{16}
\note{Pursuing Stacks, § 2 (suite ; p. 4 de l'auteur) ; imprimé dans Maltsiniotis (éd.), À la poursuite des champs, vol. I, SMF 2022. Non transcrit.}

\page{17}
\note{Pursuing Stacks, § 2 (suite ; p. 3 de l'auteur, classée après la p. 4) ; imprimé dans Maltsiniotis (éd.), À la poursuite des champs, vol. I, SMF 2022. Non transcrit.}

\page{18}
\note{Pursuing Stacks, fin du § 2 et § 3 (p. 5 de l'auteur) ; imprimé dans Maltsiniotis (éd.), À la poursuite des champs, vol. I, SMF 2022. Non transcrit.}

\page{19}
\note{Pursuing Stacks, fin du § 3 et formule finale de la lettre à Quillen (p. 5' de l'auteur) ; imprimé dans Maltsiniotis (éd.), À la poursuite des champs, vol. I, SMF 2022. Non transcrit.}

\page{20}
\note{Pursuing Stacks, § 4, daté « 20.2. » (p. 6 de l'auteur) ; imprimé dans Maltsiniotis (éd.), À la poursuite des champs, vol. I, SMF 2022. Non transcrit. Le texte se poursuit au-delà de ce lot.}

\end{document}
